\documentclass[11pt]{article}

\usepackage[margin=1in]{geometry}
\usepackage[T1]{fontenc}
\usepackage[utf8]{inputenc}
\usepackage{lmodern}
\usepackage{microtype}
\usepackage{amsmath,amssymb}
\usepackage{booktabs}
\usepackage{array}
\usepackage{tabularx}
\usepackage{enumitem}
\usepackage{hyperref}
\usepackage[numbers,sort&compress]{natbib}

\hypersetup{
  colorlinks=true,
  linkcolor=blue,
  citecolor=blue,
  urlcolor=blue
}

\title{Bounded Same-Family Latent Delegation in Gemma-2:\\Asymmetric Two-Path Routing Surpasses Bridge Baselines on Held-Out Language-Model Probes}
\author{Anonymous Submission}
\date{}

\begin{document}

\maketitle

\begin{abstract}
We study a bounded same-family latent delegation setting in which a frozen large language model keeps ownership of the master residual stream and final logits while replacing part of its middle computation with delegated computation from a frozen smaller model plus learned interface modules. We instantiate this setting with Gemma-2 9B and Gemma-2 2B under a single RTX 5090-class GPU budget and frozen backbones throughout.

A fixed contiguous delegation baseline provides only a qualified feasibility result: it improves over skip-only and no-small controls, but not over strong learned bridge baselines that remain entirely in large hidden space. A local asymmetric window search then rejects the fixed contiguous prior and yields a stable two-window shortlist. A static two-path mixture over that shortlist is the first model to beat both bridge controls on the primary internal LM-style output metrics, specifically teacher KL and held-out NLL, and a low-capacity token-wise router over the same two paths improves further.

The final token-wise model beats both bridge controls on the primary internal LM-style output metrics on a development holdout and again on an untouched confirmation holdout. Follow-up monotone-corridor and sublayer-attribution analyses strengthen the explanation of why the model works but do not produce a better model. A bounded external generalization suite shows real but mixed external validity: the cleanest carryover remains on held-out LM-style scoring, while multiple-choice gains are selective rather than broad. The strongest defensible claim is therefore a bounded same-family positive result with mixed external generalization, not a broad downstream-superiority claim.
\end{abstract}

\section{Introduction}

This paper asks a narrow question: can a large same-family model preserve the master residual stream and final logits while delegating part of its middle computation to a smaller same-family model through latent-space transfer? The framing is intentionally bounded. We do not claim thought transfer, model equivalence, or broad downstream superiority. We ask whether delegated latent computation can recover useful work under strict practical constraints.

Those constraints are deliberate. We study a single RTX 5090-class GPU setting with frozen Gemma-2 9B and Gemma-2 2B backbones, learned interface modules only, and no large-scale finetuning or distributed training. This makes the result interpretable. If the hybrid fails, it fails in a realistic frozen-backbone regime; if it succeeds, the success can be tied to a concrete delegated-computation mechanism rather than to broad adaptation capacity.

Our main finding is that the original fixed contiguous substitution is the wrong structural prior. A real-model asymmetric window search produces a near-tied two-window shortlist, and a static two-path mixture over that shortlist is the first model to beat strong large-space bridge baselines on held-out LM-style probes. A low-capacity token-wise router over the same two paths improves further and gives the best final model. The bridge win survives not only on a development holdout but also on an untouched confirmation holdout.

At the same time, the paper remains deliberately conservative. Monotone-corridor analysis and sublayer attribution strengthen the explanation of why the final model works, but neither produces a better model than the final token-wise two-path router. A bounded external generalization suite also shows that the result does not translate into broad external superiority: the clearest carryover remains on held-out LM-style scoring, while multiple-choice accuracy is mixed.

Our contributions are threefold. First, we show that fixed contiguous delegation is a poor structural prior in this same-family Gemma-2 setting, and that a bounded asymmetric window search identifies a much stronger local shortlist. Second, we show that a static two-path mixture and then a low-capacity token-wise router can beat strong bridge baselines on both a development holdout and an untouched confirmation holdout. Third, we provide explanatory follow-up analyses, monotone-corridor discovery and sublayer attribution, that strengthen the interpretation of the best model while also clarifying the limits of its external generalization.

\section{Related Work}

Our work sits between positive results on representation transfer and negative results on cross-scale incompatibility. On the positive side, recent work on same-family or closely related models suggests that hidden spaces can often be connected by simple affine or lightweight learned maps, and that transferred representations can preserve useful behavior \citep{bello2025lrt,chen2025stitching}. In language models, model stitching and related feature-transfer results similarly suggest that residual-stream features can remain compatible across models under suitable interfaces \citep{chen2025stitching}. Our design inherits that perspective, but the operational question here is stricter: we ask whether delegated small-model computation can replace a missing block of a larger model's forward pass under frozen-backbone constraints.

At the same time, recent negative results on cross-scale transfer motivate a conservative design and a strong control set \citep{tan2025incompatibility}. We therefore remain within one model family, keep all backbone weights frozen, and evaluate against learned bridge baselines that stay entirely in large hidden space. These controls matter because a positive result against weak baselines alone would not distinguish useful delegated computation from a better use of trainable capacity in the large model's own latent space.

Our contribution is not to show that same-family latent alignment is possible in principle. The contribution is operational: under frozen-backbone single-GPU constraints, fixed contiguous delegation fails, whereas asymmetric shortlist selection followed by static and token-wise two-path routing can surpass strong large-space bridges on held-out LM-style probes. The same-family choice is grounded in Gemma-2 itself \citep{gemmateam2024gemma2}. Gemma-2 2B and 9B share the same family structure while still differing enough in scale to make cross-scale mismatch meaningful. For bounded external evaluation, we use standard log-likelihood-compatible benchmarks: HellaSwag \citep{zellers2019hellaswag}, PIQA \citep{bisk2020piqa}, WinoGrande \citep{sakaguchi2020winogrande}, ARC-Easy and ARC-Challenge \citep{clark2018arc}, and LAMBADA \citep{paperno2016lambada}.

\section{Setup And Problem Statement}

We study one-way latent delegation. The large model owns the input path, the master residual stream, the suffix, and the final logits. The small model owns delegated latent computation only. Communication between the two models is through learned latent interfaces rather than token re-injection, and all backbone weights remain frozen.

The default models are Gemma-2 9B as the large model and Gemma-2 2B as the small model. The original fixed-window baseline removes large layers \texttt{24..29} and replaces them with delegated small layers \texttt{14..19}, entered from the small hidden state immediately before that delegated block. Concretely, the large prefix is layers \texttt{0..23}, the removed large block is \texttt{24..29}, the large suffix is \texttt{30..41}, the small reference hidden is taken after layer \texttt{13}, and the delegated small block is \texttt{14..19}.

The scientific question is not whether such a hybrid can merely run. The stronger question is whether delegated small-model computation adds value beyond strong controls, including \texttt{skip\_only} removal, no-small interface controls, learned bridge baselines that remain in large hidden space, and parameter-matched bridge baselines. The empirical protocol is designed to answer that stronger question directly.

\section{Method}

\subsection{Base Interface}

Let $h_t^L$ denote the large-model hidden state at token position $t$ after the frozen large prefix, and let $N(\cdot)$ denote RMSNorm. For each delegated path $p$, let $E_p$ be the entry projector into small latent space, $S_p$ the frozen delegated small-model window, and $R_p$ the return adapter back into large hidden space. The path-specific returned delta is
\begin{equation}
\Delta_{p,t} = R_p(S_p(E_p(N(h_t^L)))).
\end{equation}

For the single-path fixed-window hybrid, we write $\Delta_t := \Delta_{p,t}$ for the active delegated path. The hybrid hidden state after the removed large block is then
\begin{equation}
h_t^H = h_t^L + \Delta_t,
\end{equation}
after which the frozen large suffix produces the final logits.

\subsection{Static Two-Path Mixture}

The first successful continuation model activates two delegated paths in parallel: path B, which maps \texttt{24..27 -> 14..19}, and path A, which maps \texttt{24..27 -> 16..18}. In the static two-path mixture, the returned deltas are combined by a learned global two-logit softmax,
\begin{align}
w &= \mathrm{softmax}(\alpha),\\
\Delta_t &= w_B \Delta_{B,t} + w_A \Delta_{A,t}.
\end{align}

The matched no-small control keeps the same entry projectors, return adapters, and global mixture weights but removes actual delegated small-model computation.

\subsection{Token-Wise Two-Path Routing}

The final model replaces the global mixture with a low-capacity token-wise router. The gate reads only the large-prefix hidden state at the splice boundary. In the confirmed final configuration, the gate is a direct linear head over RMS-normalized prefix states:
\begin{align}
g_t &= \mathrm{softmax}(W_g N(h_t^L) + b_g),\\
\Delta_t &= g_{t,B}\Delta_{B,t} + g_{t,A}\Delta_{A,t}.
\end{align}

The token-wise no-small control uses the same gate family and the same interface routes but removes delegated small-model computation.

\subsection{Training Objectives}

Stage A trains only the entry projector and aligns the large-model splice state to the frozen small-model reference state:
\begin{equation}
\mathcal{L}_A = \mathrm{MSE}(E(h^L), h^S_{\mathrm{ref}}) + \mathrm{CosineLoss}(E(h^L), h^S_{\mathrm{ref}}).
\end{equation}

The decisive training regime is output-aware Stage B. Its implemented objective is
\begin{equation}
\mathcal{L}_B = \mathrm{MSE}(h^H, h^T) + \mathrm{CosineLoss}(h^H, h^T) + \lambda_{\mathrm{KL}}\mathcal{L}_{\mathrm{KL}} + \lambda_{\mathrm{CE}}\mathcal{L}_{\mathrm{CE}} + \lambda_D \lVert \Delta \rVert_2^2,
\end{equation}
where $h^T$ is the frozen large-model hidden state after the removed large block, $\mathcal{L}_{\mathrm{KL}}$ is teacher-logit KL, and $\mathcal{L}_{\mathrm{CE}}$ is shifted next-token cross-entropy. In the confirmed runs, $\lambda_{\mathrm{KL}} = 5.0$, $\lambda_{\mathrm{CE}} = 1.0$, and $\lambda_D = 10^{-4}$. The token-wise gate adds only a small stability package: a weak entropy term, a weak KL-to-static-prior term, and no temporal smoothness term in the confirmed final configuration. Stage C is intentionally not used in this paper because the main bounded claim was already established before it, while broader generalization remained mixed and an additional distillation stage would have added capacity and confounds without answering a stronger question.

\subsection{Fairness And Parameter Budgets}

Parameter matching is explicit because routing adds trainable capacity beyond a plain large-space bridge. The static two-path mixture uses \texttt{753,666} trainable parameters, the token-wise two-path router uses \texttt{764,418}, and the matched token-wise no-small control uses the same budget. The plain bridge baseline uses \texttt{458,753} trainable parameters, so we also include an updated parameter-matched bridge with \texttt{766,977} trainable parameters. This fairness audit matters because the core scientific comparison is not whether the hybrid works at all, but whether delegated small-model computation adds value beyond a strong use of comparable trainable capacity in the large model's own latent space.

\section{Experimental Protocol}

\subsection{Hardware, Seeds, And Workflow}

All experiments were run in a native Windows workflow on a single RTX 5090-class GPU using same-family Gemma-2 9B and 2B backbones. The core confirmed result families use seeds \texttt{42, 43, 44}. Deterministic evaluation subsets and saved sample IDs are used throughout.

\subsection{Development Holdout And Untouched Confirmation Holdout}

We distinguish two LM-style holdout policies. The first is a development holdout: the original held-out slice reused during model development and model-selection decisions. The second is an untouched confirmation holdout: a fresh \texttt{wikitext-103-v1} test-split slice sampled only after the winning continuation structure had been fixed. The untouched confirmation holdout contains \texttt{32} sequences at \texttt{seq\_len = 256} sampled with seed \texttt{7606}. We treat the untouched confirmation holdout as the primary basis for the strongest internal claim because it is the stricter safeguard against repeated reuse of the development slice.

\subsection{Primary Metrics}

Primary ranking is output-first: teacher KL, then NLL, then perplexity, then top-1 agreement, then top-5 overlap. KL is ranked first because the central question is whether delegated computation reproduces the functional role of the removed large-model block relative to the frozen full-large teacher. Hidden-space MSE and cosine are reported only as diagnostics.

\subsection{Bounded Generalization}

For bounded external generalization, we evaluate the frozen final model and key controls on HellaSwag, PIQA, WinoGrande, ARC-Easy, ARC-Challenge, and a held-out LAMBADA slice. Multiple-choice tasks are scored by normalized conditional answer log-likelihood; the LM-style slice is scored by KL, NLL, and perplexity. Uncertainty is reported with paired bootstrap estimates against the main internal baselines. All six external tasks use deterministic bounded subsets rather than full benchmark sweeps. HellaSwag, PIQA, WinoGrande, ARC-Easy, ARC-Challenge, and LAMBADA each use \texttt{64} examples, with fixed sampling seeds \texttt{9001}, \texttt{9002}, \texttt{9003}, \texttt{9004}, \texttt{9005}, and \texttt{9010} respectively; exact sample IDs are saved in the supplementary materials.

\section{Results}

\subsection{Fixed-Window Hybrid as a Qualified Feasibility Baseline}

The fixed-window hybrid establishes feasibility but not the final claim. In its best output-aware form, it improves over \texttt{skip\_only} removal and over the no-small interface control, showing that delegated small-model computation is real and functionally used. However, it does not beat the strong bridge baselines. The correct interpretation of this stage is therefore positive feasibility evidence with a clear structural limit.

\subsection{Asymmetric Window Search Rejects the Fixed Contiguous Prior}

A real-model local window search shows that the original fixed contiguous substitution is the wrong structural prior. The legacy \texttt{24..29 -> 14..19} candidate is substantially worse than the two candidates that later become the shortlist. After confirmation, the shortlist remains near-tied: \texttt{24..27 -> 14..19} reaches KL/NLL \texttt{0.281641 / 3.078029}, whereas \texttt{24..27 -> 16..18} reaches \texttt{0.282215 / 3.074461}. The first is slightly better on KL and top-5 overlap; the second is slightly better on NLL, perplexity, and top-1 agreement. Table~\ref{tab:progression} summarizes how that shortlist then leads to the final routing result.

\begin{table}[t]
\centering
\caption{Structural progression from the fixed-window baseline to the final routing model. Rows come from different evaluation stages and holdout policies and should not be read as a single directly comparable leaderboard.}
\label{tab:progression}
\small
\begin{tabular}{>{\raggedright\arraybackslash}p{0.20\linewidth} >{\raggedright\arraybackslash}p{0.30\linewidth} >{\raggedright\arraybackslash}p{0.20\linewidth} >{\centering\arraybackslash}p{0.08\linewidth} >{\centering\arraybackslash}p{0.07\linewidth} >{\centering\arraybackslash}p{0.07\linewidth} >{\centering\arraybackslash}p{0.08\linewidth}}
\toprule
Phase & Compared model & Holdout policy & Seeds & KL & NLL & PPL \\
\midrule
Fixed-window feasibility & Output-aware fixed-window hybrid & Development holdout & 42,43,44 & 0.655263 & 3.423486 & 30.723442 \\
Asymmetric shortlist & Best single path \texttt{24..27 -> 14..19} & Phase 1 confirmation probe & 42,43,44 & 0.281641 & 3.078029 & 21.780681 \\
Static mixture & Two-path static mixture & Development holdout & 42,43,44 & 0.267095 & 3.000438 & 20.156769 \\
Final model & Token-wise two-path routing & Development holdout & 42,43,44 & 0.255739 & 2.980182 & 19.763760 \\
\bottomrule
\end{tabular}
\end{table}

\subsection{Static Two-Path Mixture Surpasses Bridge Baselines}

The static two-path mixture is the first model to beat both bridge controls on the primary output metrics. On the development holdout, it improves over both bridge baselines in KL and NLL, and the same pattern survives on the untouched confirmation holdout. This is the first clean bridge win in the project, which is why it matters more than a simple re-ranking among single-window candidates. The matched no-small mixture also improves over the earlier no-small baselines, but the full static mixture remains stronger, indicating that the gain cannot be reduced to route mixing alone.

\subsection{Token-Wise Two-Path Routing Gives the Best Final Model}

The token-wise two-path router improves further over the static mixture and remains ahead of both bridge controls on both holdout policies. Table~\ref{tab:internal-main} gives the main comparison. On the untouched confirmation holdout, token-wise routing reaches KL/NLL \texttt{0.248886 / 3.185004}, compared with \texttt{0.267244 / 3.213048} for the static mixture, \texttt{0.289564 / 3.295081} for \texttt{bridge\_only}, and \texttt{0.301746 / 3.327024} for the parameter-matched bridge. On the development holdout, the same ordering remains: \texttt{0.255739 / 2.980182} for token-wise routing, \texttt{0.267095 / 3.000438} for the static mixture, \texttt{0.288448 / 3.072051} for \texttt{bridge\_only}, and \texttt{0.302323 / 3.102081} for the parameter-matched bridge.

The bridge comparison is clean: on both the development holdout and the untouched confirmation holdout, the token-wise model wins on the primary metrics in all three seeds against both bridge baselines. The no-small comparison is also favorable in aggregate KL and NLL on both holdouts, but it is slightly less clean on the untouched confirmation holdout, where the joint seed-level KL/NLL win count is \texttt{2/3} rather than \texttt{3/3}. That caveat is weaker than the bridge result, but it should still be stated explicitly.

The gate diagnostics are consistent with real routing rather than trivial collapse. The gate favors path A on average, retains substantial mass on path B, and remains far from full collapse. The matched no-small control uses the same gate family but is sharper, more collapsed, and materially weaker, which argues against a pure route-capacity explanation.

\begin{table}[t]
\centering
\caption{Final internal comparison on the development and untouched confirmation holdouts. Lower is better for KL and NLL.}
\label{tab:internal-main}
\small
\begin{tabular}{>{\raggedright\arraybackslash}p{0.36\linewidth} >{\centering\arraybackslash}p{0.14\linewidth} >{\centering\arraybackslash}p{0.14\linewidth} >{\centering\arraybackslash}p{0.14\linewidth} >{\centering\arraybackslash}p{0.14\linewidth}}
\toprule
Model & Dev. KL & Dev. NLL & Confirm. KL & Confirm. NLL \\
\midrule
Token-wise two-path routing & 0.255739 & 2.980182 & 0.248886 & 3.185004 \\
Token-wise no-small control & 0.257501 & 3.038605 & 0.251294 & 3.261786 \\
Static two-path mixture & 0.267095 & 3.000438 & 0.267244 & 3.213048 \\
\texttt{bridge\_only} & 0.288448 & 3.072051 & 0.289564 & 3.295081 \\
Parameter-matched bridge & 0.302323 & 3.102081 & 0.301746 & 3.327024 \\
\bottomrule
\end{tabular}
\end{table}

\subsection{Explanatory Follow-Ups}

The explanatory follow-ups strengthen interpretation rather than replacing the final model. Monotone-corridor analysis recovers a broader low-cost asymmetric region around the successful two-path shortlist, suggesting that the final model exploits a local alignment corridor rather than a lucky hard window. Sublayer attribution shows that both delegated attention and delegated MLP matter. Suppressing MLP hurts more, but suppressing attention is also materially harmful, so the gain is not well explained by a single-component simplification.

\subsection{Bounded Generalization Is Real but Mixed}

Bounded external generalization is real but mixed. The cleanest external strength remains on LM-style scoring: on the held-out LAMBADA slice, the token-wise model is favorable against both bridge baselines in KL and favorable against the parameter-matched bridge in NLL, while the NLL comparison against \texttt{bridge\_only} is positive by point estimate but not cleanly separated under paired uncertainty. The multiple-choice picture is weaker and should be stated that way. HellaSwag and ARC-Challenge are positive by point estimate, but their paired intervals cross zero; WinoGrande is mixed; and PIQA and ARC-Easy are negative relative to the bridge baselines. The strongest defensible external-validity claim is therefore that the final model is not just a single-slice artifact and retains a real LM-style held-out advantage, not that it broadly dominates strong bridge baselines across downstream tasks.

\begin{table}[t]
\centering
\caption{Bounded external generalization summary on deterministic 64-example subsets. Positive multiple-choice rows are point-estimate wins only unless paired bootstrap intervals are clearly separated; where noted in the text, some paired intervals still cross zero.}
\label{tab:generalization}
\small
\begin{tabular}{>{\raggedright\arraybackslash}p{0.16\linewidth} >{\raggedright\arraybackslash}p{0.12\linewidth} >{\centering\arraybackslash}p{0.14\linewidth} >{\centering\arraybackslash}p{0.14\linewidth} >{\centering\arraybackslash}p{0.14\linewidth} >{\centering\arraybackslash}p{0.14\linewidth}}
\toprule
Task & Metric & Token-wise & Static mix & \texttt{bridge\_only} & Param.-matched bridge \\
\midrule
HellaSwag & Accuracy & 0.671875 & 0.671875 & 0.656250 & 0.656250 \\
PIQA & Accuracy & 0.723958 & 0.723958 & 0.734375 & 0.744792 \\
WinoGrande & Accuracy & 0.645833 & 0.656250 & 0.635417 & 0.656250 \\
ARC-Easy & Accuracy & 0.791667 & 0.796875 & 0.828125 & 0.828125 \\
ARC-Challenge & Accuracy & 0.442708 & 0.427083 & 0.432292 & 0.437500 \\
LAMBADA & KL / NLL & 0.251354 / 3.423984 & 0.258699 / 3.419273 & 0.254975 / 3.433371 & 0.266066 / 3.446407 \\
\bottomrule
\end{tabular}
\end{table}

\section{Reproducibility Statement}

We release exact benchmark slice definitions, saved sample IDs, fixed seeds, canonical result tables, figure specifications, and a reproducibility manifest generated directly from frozen artifacts. Supplementary materials record commit provenance, artifact roots, and the Windows-native commands needed to rerun the reported evaluations.

\section{Limitations}

This paper reports a bounded systems-and-mechanism result, not a broad benchmark result. The strongest evidence is confined to one same-family pair, Gemma-2 9B and 2B, under a frozen-backbone single-GPU regime. The strongest external carryover remains on LM-style scoring, while broader multiple-choice generalization is mixed rather than broad. The no-small comparison is also slightly weaker on the untouched confirmation holdout than the bridge comparison. Stage C is intentionally not used because the main bounded claim was already established before it, while broader generalization remained mixed and an additional distillation stage would have added capacity and confounds without answering a stronger question. The paper does not establish cross-family robustness, broad downstream superiority, or a universal delegation principle. These limitations are part of the claim boundary, not footnotes to it.

\section{Conclusion}

The strongest result in this paper is a bounded positive one. Same-family latent delegation does not work here because a fixed hard layer match happens to be adequate. It works only after the fixed contiguous substitution is rejected, a local asymmetric shortlist is identified, and delegated computation is reformulated as a low-capacity two-path routing problem.

The resulting token-wise two-path model beats both strong bridge controls on the primary internal LM-style output metrics on a development holdout and again on an untouched confirmation holdout in the frozen Gemma-2 9B $\rightarrow$ 2B setting. This is stronger than the earlier fixed-window feasibility result and is the correct main claim of the paper. At the same time, the paper remains deliberately narrow. The explanatory follow-up branches clarify why the model works but do not replace it, and the bounded external suite shows mixed rather than broad generalization. The right conclusion is therefore narrow and strong at once: one-way same-family latent delegation can surpass strong bridge baselines inside this frozen-backbone Gemma-2 setting, but the current evidence does not justify a broad downstream-superiority claim outside that regime.

\appendix
\section{Supplementary Protocol Details}

\subsection{Bounded Generalization Subset Policy}

Table~\ref{tab:subset-policy} makes the bounded external-evaluation policy explicit. None of these tasks uses a full benchmark sweep in the current paper. Every task uses a deterministic fixed-size subset with saved sample identifiers.

\begin{table}[h]
\centering
\caption{Deterministic subset policy for the bounded external generalization suite.}
\label{tab:subset-policy}
\small
\begin{tabular}{>{\raggedright\arraybackslash}p{0.22\linewidth} >{\raggedright\arraybackslash}p{0.18\linewidth} >{\raggedright\arraybackslash}p{0.14\linewidth} >{\centering\arraybackslash}p{0.12\linewidth} >{\centering\arraybackslash}p{0.12\linewidth}}
\toprule
Task & Split & Evaluation type & Sample count & Seed \\
\midrule
HellaSwag & Validation & Multiple choice & 64 & 9001 \\
PIQA & Validation & Multiple choice & 64 & 9002 \\
WinoGrande & Validation & Multiple choice & 64 & 9003 \\
ARC-Easy & Validation & Multiple choice & 64 & 9004 \\
ARC-Challenge & Validation & Multiple choice & 64 & 9005 \\
LAMBADA & Test & LM-style scoring & 64 & 9010 \\
\bottomrule
\end{tabular}
\end{table}

\subsection{Supplementary Release Contents}

The supplementary release accompanies this manuscript with:
\begin{itemize}[leftmargin=1.5em]
\item canonical paper tables generated from frozen artifacts;
\item figure specifications and machine-readable summaries;
\item exact sample-ID files and slice definitions for the internal holdouts and bounded generalization suite;
\item a reproducibility manifest with commit provenance, artifact roots, environment metadata, and Windows-native rerun commands.
\end{itemize}

\subsection{Claim Boundary Reminder}

The supported claim remains bounded. The strongest evidence is internal to the frozen Gemma-2 9B $\rightarrow$ 2B setting and is anchored on LM-style output metrics. Broader multiple-choice generalization remains mixed, and Stage C is intentionally excluded from the present paper.


\bibliographystyle{plainnat}
\bibliography{references}

\end{document}
